% !TEX program = lualatex
\documentclass[aspectratio=169,professionalfonts]{beamer}
\newcommand{\slidemode}{1}  % カラー投影モード
\usepackage{beamer_template}
\usepackage{enumerate}

\newcommand{\LectureNum}{2}
\newcommand{\LectureTitle}{相似性と移動法則}
\newcommand{\CourseName}{応用数学}
\renewcommand{\TermName}{前期}

\begin{document}

\begin{frame}{第2回　相似性と移動法則}
  \begin{block}{本回の内容}
  相似性と移動法則
  \end{block}
  \vfill\hfill{\scriptsize \textcolor{gray}{第2章 ラプラス変換　§1}}
\end{frame}

\begin{frame}{定理：相似性}
  \begin{block}{}
  \[
    L[f(at)] = (1/a) F(s/a)  (a は正の定数)
  \]
  \end{block}
  \vfill\hfill{\scriptsize \textcolor{gray}{p.\,50--53}}
\end{frame}

\begin{frame}{定理：像関数の移動法則}
  \begin{block}{}
  \[
    L[e^{at}f(t)] = F(s-a)  (a は定数)
  \]
  \end{block}
  \vfill\hfill{\scriptsize \textcolor{gray}{p.\,50--53}}
\end{frame}

\begin{frame}{定理：原関数の移動法則（ヘビサイドの移動定理）}
  \begin{block}{}
  \[
    L[f(t-\mu)U(t-\mu)] = e^{-\mu s}F(s)  (\mu \geq 0)
  \]
  \[
    L^{-1}[e^{-\mu s}F(s)] = f(t-\mu)U(t-\mu)
  \]
  \end{block}
  \vfill\hfill{\scriptsize \textcolor{gray}{p.\,50--53}}
\end{frame}

\begin{frame}{例題7}
  \begin{exampleblock}{問題}
    $L[\sin(\omega t)]$ を求めよ（ $\omega$ は 0 でない定数）\\[2pt]
  \end{exampleblock}
\end{frame}

\begin{frame}{例題7　解答}
  \begin{exampleblock}{解答}
  \begin{enumerate}[1.]
    \item[] $\omega > 0$ のとき : 相似性と例題 4 の結果から
    \item[] $L[\sin(\omega t)] = (1/\omega) \cdot 1/((s/\omega)^{2}+1) = \omega/(s^{2}+\omega^{2})$
    \item[] $\omega < 0$ のとき $: \sin(-\omega t) = -\sin(\omega t)$ を利用
    \item[] $L[\sin(\omega t)] = -L[\sin(-\omega t)] = -(-\omega)/(s^{2}+\omega^{2}) = \omega/(s^{2}+\omega^{2})$
  \end{enumerate}
  \end{exampleblock}
  \begin{alertblock}{答}
    $L[\sin(\omega t)] = \omega/(s^{2}+\omega^{2})$（$\omega$ の正負に関係なく）
  \end{alertblock}
\end{frame}

\begin{frame}{例題8}
  \begin{exampleblock}{問題}
    $L[e^{\alpha t}\sin(\mu t)]$ を求めよ（ $\alpha$ は定数、 $\mu \neq 0$ ）\\[2pt]
  \end{exampleblock}
\end{frame}

\begin{frame}{例題8　解答}
  \begin{exampleblock}{解答}
  \begin{enumerate}[1.]
    \item[] 例題 7 より $L[\sin(\mu t)] = \mu/(s^{2}+\mu^{2}) (s > 0)$
    \item[] 像関数の移動法則 $L[e^{\alpha t}f(t)] = F(s-\alpha)$ を適用すると
    \item[] s を $s-\alpha$ に置き換えて
    \item[] $L[e^{\alpha t}\sin(\mu t)] = \mu/((s-\alpha)^{2}+\mu^{2})$
  \end{enumerate}
  \end{exampleblock}
  \begin{alertblock}{答}
    $L[e^{\alpha t}\sin(\mu t)] = \mu/((s-\alpha)^{2}+\mu^{2})$
  \end{alertblock}
\end{frame}

\begin{frame}{演習問題（p.\,50--53）}
  \begin{exampleblock}{自分で解いてみよう}
  \begin{description}
    \item[問8] $L[\cos(\omega t)]$ を求めよ（ $\omega$ は 0 でない定数）
    \item[問9] $\cos^{2}t, \sin^{2}t, \sin t\cdot \cos t$ のラプラス変換を求めよ（平方の公式を使用）
    \item[問10] $L[e^{\mu t}\cdot \cos(\mu t)]$ を求めよ（ $\mu$ は定数）
    \item[問11] $L[\sin(t - \pi/4)\cdot U(t - \pi/4)]$ を求めよ
    \item[問12] 関数 $f(t) = (t-2)U(t-2)$ のグラフをかけ。また $L[f(t)]$ を求めよ
  \end{description}
  \end{exampleblock}
\end{frame}

\end{document}
